\documentclass[11pt]{article}
\usepackage[margin=0.75in]{geometry}
\usepackage{hyperref}
\usepackage{amsmath}
\usepackage{setspace}

\hypersetup{colorlinks=true, urlcolor=blue!60!black, citecolor=blue!60!black}
\setstretch{1.1}

\title{\vspace{-1.5cm}\normalsize\textbf{Final Project Proposal:\\
Prior Sensitivity in Bayesian Logistic Regression:\\
A Small-Sample Medical Study}}
\author{
    William Chen\\
    \textit{University of California, Irvine}\\
    Instructor: Prof.\ Weining Shen\\
    \textit{Stat 205P: Bayesian Data Analysis}
}
\date{}

\begin{document}
\maketitle
\thispagestyle{empty}

\section*{Introduction}

This project applies Bayesian logistic regression to a small medical
dataset to examine how prior choice shapes posterior inference.
A fundamental property of Bayesian inference is that the influence
of the prior diminishes as sample size grows. In small clinical
studies, common in medicine due to cost and ethical constraints,
prior choice can substantially alter conclusions. Understanding
this dependence is critical for practitioners who must choose
priors carefully and communicate uncertainty honestly.

\section*{Dataset}

I will use the \texttt{birthwt} dataset ($n = 189$) from R's
\texttt{MASS} package, collected at Baystate Medical Center,
Springfield, MA (1986). The binary outcome is low birth weight
($< 2.5$ kg), and predictors include maternal age, weight, race,
smoking status during pregnancy, history of previous premature
labour (\texttt{ptd}), history of hypertension, uterine irritability,
and number of physician visits during the first trimester.

\section*{Methods}

Two analyses are planned:

\begin{enumerate}
    \item \textbf{Prior Sensitivity Analysis} (full data, $n = 189$):
    Fit Bayesian logistic regression under three priors: diffuse
    $\mathcal{N}(0, 100)$, weakly informative $\mathcal{N}(0, 2.5)$,
    and clinically informative priors motivated by medical literature
    (e.g., positive coefficients for smoking and hypertension).
    Compare posterior means and 95\% credible intervals across models.

    \item \textbf{Subsampling Experiment}: Repeatedly subsample at
    $n \in \{20, 40, 80, 189\}$ and refit all three models. Visualize
    how posterior distributions and credible interval widths evolve
    with sample size, illustrating the transition from
    prior-dominated to likelihood-dominated inference.
\end{enumerate}

All analyses will be conducted in R using \texttt{rstanarm} and
visualized with \texttt{bayesplot} and \texttt{ggplot2}.

\section*{Goals}

\begin{itemize}
    \item Demonstrate how prior sensitivity varies with sample size
    in a real medical context.
    \item Identify which risk factors (e.g., smoking, hypertension)
    are most affected by prior choice under small $n$.
    \item Provide practical guidance on prior selection for small
    clinical datasets, contrasting diffuse, weakly informative,
    and domain-informed priors.
\end{itemize}

\end{document}